\subsection{Binary Tool Usage}
\subsubsection*{Quick Start}
Generate only the intermediate Python script (\texttt{.v.py}):
\begin{verbatim}
pytv path/to/design.pytv
\end{verbatim}

Generate Verilog/\texttt{.inst} and keep \texttt{.v.py}:
\begin{verbatim}
pytv path/to/design.pytv --run-py
\end{verbatim}

Generate Verilog/\texttt{.inst} and delete \texttt{.v.py}:
\begin{verbatim}
pytv path/to/design.pytv --run-py-del
\end{verbatim}

Set an explicit output base file:
\begin{verbatim}
pytv path/to/design.pytv -o build/design.v -R
\end{verbatim}

\subsubsection*{Examples in This Repository}
\begin{verbatim}
cargo run -- examples/test.pytv -R
cargo run -- examples/D1.pytv -R \
  --preamble examples/D1_preamble.py
\end{verbatim}

\subsubsection*{Command-Line Options}
Command-line options are summarized in Table~\ref{tab:cli-options}.
\begin{table}[htbp]
\centering
\small
\caption{PyTV command-line options.}
\label{tab:cli-options}
\begin{tabularx}{\linewidth}{@{} l X @{}}
\toprule
\textbf{Option} & \textbf{Description} \\
\midrule
\texttt{INPUT} & Input \texttt{.pytv} file (required positional argument). \\
\texttt{-o, --output FILE} & Set output Verilog path/base. Derived \texttt{.v.py} and \texttt{.inst} names follow this base. \\
\texttt{-r, --run-py} & Execute generated Python and keep \texttt{.v.py}. \\
\texttt{-R, --run-py-del} & Execute generated Python and delete \texttt{.v.py} after success. \\
\texttt{-t, --tab-size INT} & Tab width used when normalizing tab characters in input lines (default: 4). \\
\texttt{-m, --magic STRING} & Magic marker string used after \texttt{//} and \texttt{/*} (default: \texttt{!}). \\
\texttt{-v, --var KEY=VAL} & Inject Python variables; option can be repeated. \\
\texttt{-p, --preamble FILE} & Prepend a Python preamble script before template conversion. \\
\bottomrule
\end{tabularx}
\normalsize
\end{table}

\subsection{Rust Library Usage}
PyTV can be integrated directly in Rust applications.
The core types are \texttt{Config}, \texttt{FileOptions}, and \texttt{Convert}.

\subsubsection*{Programmatic Conversion Example}
\begin{verbatim}
use pytv::{Config, Convert, FileOptions};

fn main() -> Result<(), Box<dyn std::error::Error>> {
    let config = Config::new(
        "!".to_string(),
        Config::default_template_re(),
        true,  // run generated Python
        true,  // delete .v.py after successful run
        4,     // tab size
    );

    let files = FileOptions {
        input: "examples/test.pytv".to_string(),
        output: Some("build/test.v".to_string()),
    };

    let vars = Some(vec![
        ("if_rst".to_string(), "True".to_string()),
        ("if_en".to_string(), "False".to_string()),
    ]);

    let convert = Convert::new(config, files, vars, None);
    convert.convert_to_file()?;
    Ok(())
}
\end{verbatim}

\subsection{Python Binding Usage}
For Python API usage, execution modes, context behavior, and error model,
see Section~\ref{sec:tverilog}.

\subsection{VS Code Extension}
A VS Code extension with syntax highlighting is available for PyTV.
Install by searching for \texttt{pytv} in the VS Code Marketplace,
or visit the
\href{https://marketplace.visualstudio.com/items?itemName=Teddy-van-Jerry.pytv}{VS Code Marketplace page}.
